%%%%%%%%%%%%%%%%%%%%%%%%%%%%%%%%%%%%%%%%%%%%%%%%%%%%%%%%%%%%%%%%%%%%%%%%%%%%%%%
%%
%%  MARK: Packages declarations
%%
%%%%%%%%%%%%%%%%%%%%%%%%%%%%%%%%%%%%%%%%%%%%%%%%%%%%%%%%%%%%%%%%%%%%%%%%%%%%%%%

%% AMS package
\usepackage{amssymb} %% AMS Symbols
\usepackage{amsmath} %% AMS Math Macros

%% Graphic package
\usepackage{graphicx} %% Include Drawings
\usepackage[innercaption]{sidecap} %% For caption on the side of the figure

%% Increase length page package
\usepackage{needspace} %% Enhance page space

%% Smart fonts for \mathcal replacement
%\usepackage{bickham}
%\usepackage{boondox-cal}
%\usepackage{dutchcal}
\usepackage[mathcal,mathbf,text-hat-accent]{euler} %% For Math Script

%ù Add PDF page package
\usepackage{pdfpages} %% Include PDF pages

%% tikz Package
\usepackage{tikz-cd} %% Include Diagrams  http://texdoc.net/texmf-dist/doc/latex/tikz-cd/tikz-cd-doc.pdf
\usetikzlibrary{calc}
\usetikzlibrary{matrix}

\tikzcdset{
arrow style=tikz,
diagrams={>={Straight Barb[scale=0.8]}}
}

%% Include Dingbat fonts
\usepackage{pifont}

% For strike out
\usepackage[normalem]{ulem}

% For inteernet reference
\usepackage{url}

% For book name in Chinese
\usepackage{CJKutf8}
\usepackage[utf8]{inputenc} % optional

% To underline cleanly letters that go under the base line
\usepackage{contour}
\setlength{\ULdepth}{1.8pt}
\contourlength{0.8pt}

\newcommand{\smuline}[1]{
\uline{\phantom{#1}}%
\llap{\contour{white}{#1}}%
}

%%%%%%%%%%%%%%%%%%%%%%%%%%%%%%%%%%%%%%%%%%%%%%%%%%%%%%%%%%
%%
%% MARK: Abstract
%%
%%%%%%%%%%%%%%%%%%%%%%%%%%%%%%%%%%%%%%%%%%%%%%%%%%%%%%%%%%

%% Redefinition of Abstract only for Lectures on Diffeology
\renewenvironment{abstract}[1]{\noindent\rule{\linewidth}{0.4mm}\vspace{1mm}
\newline\noindent\ #1}
{\newline\noindent\rule{\linewidth}{0.4mm}\bigskip}

%%%%%%%%%%%%%%%%%%%%%%%%%%%%%%%%%%%%%%%%%%%%%%%%%%%%%%%%%%%%%%%%%%%%%%%%%%%%%%%
%%
%% MARK: Dingbat characters
%%
%%%%%%%%%%%%%%%%%%%%%%%%%%%%%%%%%%%%%%%%%%%%%%%%%%%%%%%%%%%%%%%%%%%%%%%%%%%%%%%

\def\pencil{\raisebox{-.5pt}{\Large\ding{46}}} %% A pencil picture to mark exercise

\def\nib{\raisebox{-2.5pt}{\huge\ding{49}}} %% A nib of a pen to mark the solution of an exercise

\def\blacksnow{\raisebox{-.5pt}{\ding{81}}} %% A bold black asterisc

%%%%%%%%%%%%%%%%%%%%%%%%%%%%%%%%%%%%%%%%%%%%%%%%%%%%%%%%%%
%%
%% MARK: Paragraphs: Article, Note, Question
%%
%%%%%%%%%%%%%%%%%%%%%%%%%%%%%%%%%%%%%%%%%%%%%%%%%%%%%%%%%%
\newcounter{art}
\setcounter{art}{0}

\newenvironment{Article}[2]{\refstepcounter{art}\par\noindent{\textbf{\smuline{\theart.~#1}}} #2}{}

\renewenvironment{proof}{\noindent \raisebox{-1.5pt}{\nib} {Proof.}} {\nolinebreak $\blacktriangleright$}

\chardef\atcode=\catcode`\@  % save catcode of at sign
\catcode `\@ = 11

\let\bgroup={%
\newcommand{\Note}{\@ifnextchar\bgroup\@numnote\@nonumnote}
\def\@numnote#1{\ifnum\lastskip=0\vskip .15\baselineskip\fi\noindent\uline{Note #1}}
\def\@nonumnote{\ifnum\lastskip=0\vskip .15\baselineskip\fi\noindent\uline{Note.}\ }

\catcode`\@=\atcode    % restore original catcode of at sign

\def \Question #1{\vskip-\lastskip\vspace{.25\baselineskip plus 2pt minus 2pt} \noindent{\bf #1}}

%%%%%%%%%%%%%%%%%%%%%%%%%%%%%%%%%%%%%%%%%%%%%%%%%%%%%%%%%%%%%%%%%%%%%%%%%%%%%%%
%%
%% MARK: Exercices Environment Implementation
%%
%%%%%%%%%%%%%%%%%%%%%%%%%%%%%%%%%%%%%%%%%%%%%%%%%%%%%%%%%%%%%%%%%%%%%%%%%%%%%%%

%% Exercices Environment
\newtheoremstyle{exercise} %  ⟨name⟩
{7pt}                      %  ⟨Space above⟩
{7pt}                      %  ⟨Space below⟩
{}                         %  ⟨Body font⟩
{0pt}                      %  ⟨Indent amount⟩
{}                         %  ⟨Theorem head font⟩
{.\ }                      %  ⟨Punctuation after theorem head⟩
{0pt}                      %  ⟨Space after theorem head⟩
{}                         %  ⟨Theorem head spec (can be left empty, meaning ‘normal’)⟩

\newtheorem*{exercise}{\mbox{\pencil} {\sc Exercise}}

%%%%%%%%%%%%%%%%%%%%%%%%%%%%%%%%%%%%%%%%%%%%%%%%%%%%%%%%%%
%%
%% MARK: Page layout
%%
%%%%%%%%%%%%%%%%%%%%%%%%%%%%%%%%%%%%%%%%%%%%%%%%%%%%%%%%%%

\parindent 0mm
\parskip .5ex plus 2pt

%% For Beijing WPC Lectures on Diffeology
\paperheight 225mm
\paperwidth 170mm

\marginparwidth = 0pt

%% For Beijing WPC Lectures on Diffeology
\setlength{\textheight}{180mm}
\setlength{\textwidth}{120mm}

%% For Beijing WPC
%\ usepackage[width=170mm, height=225mm, cam, noinfo, center]{crop} %% Not used now

\calclayout

\setlength{\parindent}{0em}
\setlength{\parskip}{1ex}

%%%%%%%%%%%%%%%%%%%%%%%%%%%%%%%%%%%%%%%%%%%%%%%%%%%%%%%%%%%%%%%%%%%%%%%%%%%%%%%
%%
%% MARK: Font encoding and display
%%
%%%%%%%%%%%%%%%%%%%%%%%%%%%%%%%%%%%%%%%%%%%%%%%%%%%%%%%%%%%%%%%%%%%%%%%%%%%%%%%

%% This font configuration for keeping the main text in the typewriter/draft style

\usepackage[T1]{fontenc}

\usepackage[activate]{microtype}
\usepackage[supspaced=.04em]{superiors}

\usepackage[variablett,nomath]{lmodern}
\renewcommand{\familydefault}{\ttdefault}
\usepackage[LGRgreek,frenchmath,symbolmisc]{mathastext} %% https://jfbu.github.io/mathastext/showcase.html
\MTgreekfont{lmtt}
\Mathastext[mathtt]

%%%%%%%%%%%%%%%%%%%%%%%%%%%%%%%%%%%%%%%%%%%%%%%%%%%%%%%%%%%%%%%%%%%%%%%%%%%%%%%
%% Font for Enxin Wu's Appendix
%% The Appendix must be in Charter (bch). Placing this configuration here,
%% after all other font settings, is the only reliable way to achieve this
%% without breaking the main document's font.
%%%%%%%%%%%%%%%%%%%%%%%%%%%%%%%%%%%%%%%%%%%%%%%%%%%%%%%%%%%%%%%%%%%%%%%%%%%%%%%
\usepackage{palatino} % For Enxin Appendix
\renewcommand{\familydefault}{bch} % ppl bch ugm put
\MTgreekfont{eulergreek}
\Mathastext[palatino]

%%%%%%%%%%%%%%%%%%%%%%%%%%%%%%%%%%%%%%%%%%%%%%%%%%%%%%%%%%%%%%%%%%%%%%%%%%%%%%%
%%% Redefinition of letters and symbols
%%%%%%%%%%%%%%%%%%%%%%%%%%%%%%%%%%%%%%%%%%%%%%%%%%%%%%%%%%%%%%%%%%%%%%%%%%%%%%%

\def\infty{\inftypsy}
\DeclareMathSymbol{\inftyreg}{\mathord}{symbols}{49}

\def \epsilon{\varepsilon}

\DeclareMathSymbol{\pphi}{\mathord}{mtpsymbol}{"66}
\def \varphi{\phi}

\linespread{1.08}

%%% Look for regulare expression \\begin\{article\}\{.*?\.\} beginning wih '\begin{Article}{' and ending with '.}'

%%%%%%%%%%%%%%%%%%%%%%%%%%%%%%%%%%%%%%%%%%%%%%%%%%%%%%%%%%%%%%%%%%%%%%%%%%%%%%%
%%
%% MARK: General Display
%%
%%%%%%%%%%%%%%%%%%%%%%%%%%%%%%%%%%%%%%%%%%%%%%%%%%%%%%%%%%%%%%%%%%%%%%%%%%%%%%%

\DeclareRobustCommand{\SkipTocEntry}[4]{}

\allowdisplaybreaks
